% !TEX TS-program = xelatex
% !TEX encoding = UTF-8 Unicode
% !Mode:: "TeX:UTF-8"

\documentclass{resume}
\usepackage{zh_CN-Adobefonts_external} % Simplified Chinese Support using external fonts (./fonts/zh_CN-Adobe/)
%\usepackage{zh_CN-Adobefonts_internal} % Simplified Chinese Support using system fonts
\usepackage{linespacing_fix} % disable extra space before next section
\usepackage{cite}

\begin{document}
\pagenumbering{gobble} % suppress displaying page number

\name{汪星鹏}

% {E-mail}{mobilephone}{homepage}
% be careful of _ in emaill address
\contactInfo{(+86) 18455656782}{22332102@zju.edu.cn}{}{GitHub @deeplerming}
% {E-mail}{mobilephone}
% keep the last empty braces!
%\contactInfo{xxx@yuanbin.me}{(+86) 131-221-87xxx}{}
 
% \section{个人总结}
% 本人在校成绩优秀、乐观向上，工作负责、自我驱动力强、热爱尝试新事物，认同开放、连接、共享的Web在未来的不可替代性。在校期间长期从事可视分析(Web的2D/3D时空可视化)相关研究，对Web技术发展趋势及前端工程化解决方案有浓厚兴趣。\textbf{现任职于 BAT 集团。}

% \section{\faGraduationCap\ 教育背景}
\section{教育背景}
\datedsubsection{\textbf{浙江大学},控制科学与工程学院,\textit{Fast-Lab 硕士研究生}}{2023.9 - 2026.6}
\ {研究方向：SLAM GNSS等多传感器融合定位}
\datedsubsection{\textbf{江南大学},机械工程学院,\textit{工学学士}}{2019.9 - 2023.6}
{毕业设计：基于语义分割的芯片缺陷检测系统设计}

\section{科研经历}
\begin{itemize}
	\item {\textbf{基于环境感知的空地两栖机器人鲁棒定位系统}}\hfill {2023--2024}
	\begin{itemize} 
		\item {提出了一种专门为基于被动轮的设计的环境感知鲁棒定位系统，该系统具有两个被动轮和一个标准四旋翼飞行器。该定位系统紧密集成了来自多个传感器的数据，包括摄像头、惯性测量单元 (IMU)、编码器和单点激光距离传感器。}
		\item {提出了一种异常感知方法，可以感知异常传感器并相应地动态调整优化权重。通过明确估计环境条件，例如地面地形坡度和视觉信息质量，机器人可以在地面上获得准确且鲁棒的定位结果。}
		\item {为了评估我们的工作，我们收集了各种环境下的 15 个数据集，包括平坦地形、不平坦地形、不同运动模式、相机遮挡以及起飞和着陆期间，在各种具有挑战性的场景中进行了广泛的实验，证明了我们的系统在实际应用中的有效性和可靠性。}
	\end{itemize}
	\item {\textbf{基于SLAM系统的Tof相机深度补全}} \hfill {2023--2025} 
	\begin{itemize}
		\item {Time-of-Flight (ToF)相机因其轻量性和高精度而被广泛用于机器人任务，然而其深度视距受限问题使其目前只适合小型室内场景。本工作旨在通过深度补全扩展ToF相机的深度范围，以提升其在大规模室内外场景的下游任务表现。}
		\item {在移动重建而非传统静止累积的范式下，构建了首个基于激光三维重建的大规模场景ToF深度补全数据集及基准测试，在深度范围/填充密度/采集效率等方面相较现有数据集具有显著优势。通过引入新型3D分支设计，3D-2D联合传播池化模块，和多模态跨协方差注意力，设计了针对性的深度补全网络架构。该架构可进一步利用RGB-ToF序列的视觉SLAM跟踪点云，获取远近一致引导深度，以提高深度补全质量。对模型进行轻量化适配，使其直接在无人机机载电脑实时运行，有效助力建图与自主规划等机器人下游任务。}
		\item {定量实验结果展示了相较于现有SOTA深度补全方法的卓越精度（平均绝对误差降低8.2\%以上），实现将ToF相机的有效视距拓展到数倍乃至数十几倍。两项实机下游任务实验证明，所提出方法显著提高了无人机建图任务的效率和地图完整性，提升了无人机规划器的安全性与轨迹最优性。}
	\end{itemize}
	\item {\textbf{不依赖固定GNSS基站的多机器人高精度集群定位系统}} \hfill {2024--至今}
	\begin{itemize}
		\item {传统GNSS受限环境下定位精度下降，高精度定位依赖固定基站或网络基站，通过SLAM系统解除基站固定的限制从而实现多机器人协同定位算法，同时保证鲁棒性与固定基站RTK的厘米级精度。}
		\item 
	\end{itemize}
	\end{itemize}

% \section{\faCogs\ IT 技能}


% \end{itemize}

\section{项目经历}

\datedsubsection{\textbf{紧耦合GNSS/LiDAR/IMU的智能车辆连续、无漂移状态估计器}}{2023-2024}
\begin{itemize}
	\item {针对智能车辆在复杂环境下的连续定位时里程计的累计漂移问题，设计了一种紧耦合GNSS/LiDAR/IMU的状态估计器，通过GNSS原始数据观测，减少系统定位累计误差，同时定位精度达到厘米级。}
	\item {设计了一种基于图优化的多传感器融合框架，将GNSS伪距、载波相位观测值与LiDAR里程计、IMU预积分观测值进行紧耦合优化，实现高精度、低漂移的状态估计。}
	\item {通过在多种复杂环境下进行实车测试，验证了所提出方法在连续定位中的有效性和鲁棒性，显著提升了智能车辆的导航性能。}
\end{itemize}
\datedsubsection{\textbf{林业测绘环境中的高精度回环定位系统}}{2023-2024}
\begin{itemize}
	\item {林业测绘过程中，由于树木的高度和密度变化较大，基于视觉或LiDAR的回环检测方法在该场景下表现不佳，设计了一种基于树木几何关系的高精度回环定位系统。}
	\item {利用柱状特征确定局部地图中树的位置，再由稳定且能保证旋转不变性的三角描述子来锚定每一帧观测的树间几何关系，从而快速确定是否触发回环，约束状态估计。}
	\item {通过在多种复杂环境下进行实机测试，验证了该定位器在林业测绘中的有效性和鲁棒性，在尽量减少占用计算资源的前提下提高了系统的定位精度。}
\end{itemize}

\datedsubsection{\textbf{利用无人机第三视角简化履带车控制}}{2025.10-至今}
\begin{itemize}
  \item \textbf{利用海量用户定位数据，对城市空间及人群移动特征进行研究。}第一个课题是基于香农熵和人群出行模式，构建城市网格与用户矩阵分析城市多样性/流动性分布；可视分析平台前端与可视化基于D3/Vue/Express开发，数据分析与存储采用Python/MySQL/MongoDB技术，为了均衡大数据情况下的页面可视化渲染消耗用canvas替代svg。第二个课题是对海量商场定位数据做人群分类与可视化查询，依据该系统撰写的论文被CIKM 2016(DAVA Workshop)录用，并收录于中科院软件所年会成果集
  \item 负责数据科学部HQ LAB的可视化原型开发，主导 TalkingMind 平台系统设计与前端开发
\end{itemize}

% \begin{onehalfspacing}
% \end{onehalfspacing}

% \datedsubsection{\textbf{DID-ACTE} 荷兰莱顿}{2015年}
% \role{本科毕业设计}{LIACS 交换生}
% 利用结巴分词对中国古文进行分词与词性标注，用已有领域知识训练形成 classifier 并对结果进行调优
% \begin{onehalfspacing}
% \begin{itemize}
%   \item 利用结巴分词对中国古文进行分词与词性标注
%   \item 利用已有领域知识训练形成 classifier, 并用分词结果进行测试反馈
%   \item 尝试不同规则，对 classifier 进行调优
% \end{itemize}
% \end{onehalfspacing}

\section{论文与专利}
% increase linespacing [parsep=0.5ex]
\begin{itemize}[parsep=0.2ex]
%   \item LeetCodeOJ Solutions, \textit{https://github.com/hijiangtao/LeetCodeOJ}
  \item {He W, \textbf{Wang X}, Wang P, et al. EAR-SLAM: Environment-Aware Robust Localization System for Terrestrial-Aerial Bimodal Vehicles}, ICRA2025
%   \item 中国机器人大赛暨Robocup公开赛(武术擂台赛)全国一等奖,2013年10月
  \item {Juncheng Chen, Tiancheng Lai, \textbf{Xingpeng Wang}, Bingxin Liao, Baozhe Zhang, Chao Xu*, and Yanjun Cao* ToFormer: Large-Scale Scenario Depth Completion for Time-of-Flight Camera},	T-ASE在投
%   \item 电视节目"爸爸去哪儿"可视化分析展示, \textit{https://hijiangtao.github.io/variety-show-hot-spot-vis/}
\end{itemize}

\section{技术能力}
% increase linespacing [parsep=0.5ex]
\begin{itemize}[parsep=0.2ex]
  \item \textbf{编程语言}: JavaScript (ECMAScript, Node.js), HTML/CSS, Python, Go, SQL, C, Shell
  \item \textbf{操作系统,数据库与工程构建}: Linux/macOS/MySQL/MongoDB/Git/webpack/Progressive Web App
  \item \textbf{关键词}: React/Vue.js/D3.js(SVG)/three.js(canvas, WebGL)/chrome extension/Express
\end{itemize}

% \section{\faHeartO\ 项目/作品摘要}
% \section{项目/作品摘要}
% \datedline{\textit{An Integrated Version of Security Monitor Vis System}, https://hijiangtao.github.io/ss-vis-component/ }{}
% \datedline{\textit{Dark-Tech}, https://github.com/hijiangtao/dark-tech/ }{}
% \datedline{\textit{融合社交网络数据挖掘的电视节目可视化分析系统}, https://hijiangtao.github.io/variety-show-hot-spot-vis/}{}
% \datedline{\textit{LeetCodeOJ Solutions}, https://github.com/hijiangtao/LeetCodeOJ}{}
% \datedline{\textit{Info-Vis}, https://github.com/ISCAS-VIS/infovis-ucas}{}


% \section{\faInfo\ 社会实践/其他}
% \section{社区参与/实践其他}
% % increase linespacing [parsep=0.5ex]
% \begin{itemize}[parsep=0.2ex]
%   \item 乐于参与开源社区讨论,\textbf{参与翻译 Vue.js, webpack, WebAssembly, Babel 文档，印记中文成员}
%   \item 中国科学院大学2016秋季学期可视化与可视分析课程助教，\textit{http://vis.ios.ac.cn/infovis-ucas/}
%   \item 未来论坛学生会成员、北理社联新闻信息中心主任、北理工软件学院学生会宣传部副部长(2012-2016)
%   \item 2013-2015 北京市共青团“温暖衣冬”志愿者,第九届园博会志愿者,2014 FLL机器人世锦赛志愿者
% \end{itemize}

%% Reference
%\newpage
%\bibliographystyle{IEEETran}
%\bibliography{mycite}
\end{document}
